\SourcePage{226}
\subsection{双纤维化}\label{sec:15e}
现在考察形如下图的结构：
\SourceEquationDiagram{diagram-15-11}{0.2388459}{15.11}
这里 $\Psi$ 为 Lefschetz 纤维化，总空间维数为 $2n+4\geq6$。中间空间 $\mathcal S$ 为带角正合辛四维流形，$\varpi$ 与 $\psi$ 均为带奇点的纤维化。假设 $\psi$ 没有临界点，且其纤维为带边界的连通曲面；特别地，$\mathcal S$ 只有余维 $\leq2$ 的角。对 $\varpi$，只要求它在 $\partial^h\mathcal S$ 上没有临界值，此处这部分边界是关于 $\psi$ 定义的。\znote{218}{“双”纤维化指两个纤维化叠在一起：$\varpi:\mathcal E\to\mathcal S$ 与 $\psi:\mathcal S\to\mathcal W$ 都以曲面为底，复合 $\Psi=\psi\circ\varpi$ 是 Lefschetz 纤维化。固定 $w$ 后，$\mathcal E_w\to\mathcal S_w$ 本身又是以曲面为底的纤维化；当 $w$ 取 $\Psi$ 的临界值时，$\mathcal E_w$ 上恰有一个临界点，这一纤维化随之退化。这样，关于 $2n+4$ 维总空间的问题便可化归为曲面 $\mathcal S_w$ 上的曲线配置：消失圈与匹配路径都生活在曲面里，而 $\Psi$ 的单值变换由 $\operatorname{Critv}(\varpi)$ 在曲面间的运动描述。[来源：原书 §(15e)，印刷页 218。]}取点 $w\in\mathcal W$，将 $\varpi$ 限制到该点上方的纤维，记所得映射为
\begin{equation}\label{eq:15.12}
\mathcal E_w\xrightarrow{\varpi_w}\mathcal S_w.
\end{equation}
除非 $w$ 是 $\Psi$ 的临界值，否则它就是带奇点的纤维化；其余必要性质均由关于 $\varpi,\psi,\Psi$ 的假设推出。因此，可将 $\mathcal E\to\mathcal W$ 视为一个 Lefschetz 纤维化，其光滑纤维又在曲面 $\mathcal S_w$ 上纤维化。若还满足下述附加条件，就称\eref{15.11}为\emph{（正合）Lefschetz 双纤维化}：
\begin{itemize}
\item $\varpi$ 的临界点一般，即 $D\varpi$ 在 $\Hom_{\CC}(T\mathcal E,\varpi^*T\mathcal S)$ 中，与零截面以及复秩为一的线性映射所成的层均横截。若 $x$ 是 $\varpi$ 的临界点但为 $\Psi$ 的正则点，则 $I_{\mathcal E}$ 在纤维 $\mathcal E_w$，$w=\Psi(x)$，上的限制在 $x$ 附近局部可积。此外，对每个 $w\in\partial\mathcal W$，\eref{15.12}均为 Lefschetz 纤维化。
\item 取 $w\in\operatorname{Critv}(\Psi)$，令 $x$ 为 $\mathcal E_w$ 中唯一临界点。要求 $D^2\Psi_x$ 在 $\ker(D\varpi_x)\subset T\mathcal E_x$ 上的限制非退化。殆复结构 $I_{\mathcal S}$ 在 $\varpi(x)$ 的某个邻域中可积。此外，$\mathcal E_w\setminus\{x\}$ 中 $\varpi_w$ 的所有临界点均须非退化，其 $\varpi$ 值两两不同，且均不同于 $\varpi(x)$。
\end{itemize}
下面逐步解释这些条件的含义。一般性与维数计算表明，$D\varpi$ 从不为零，其复秩为一的点所成的集合 $\operatorname{Crit}(\varpi)$ 是光滑二维子流形。事实上，

\StatementNumber{6}\begin{lemma}\label{st:15.6}
$\operatorname{Crit}(\varpi)$ 是光滑的 $I_{\mathcal E}$ 全纯曲线。
\end{lemma}
\begin{proof}
对 $x\in\operatorname{Crit}(\varpi)$ 及 $X\in T\mathcal E_x$，考虑部分二阶导数
\begin{equation}\label{eq:15.13}
\bar D^2\varpi_{x,X}:\ker(D\varpi_x)\longrightarrow\operatorname{coker}(D\varpi_x).
\end{equation}
$\operatorname{Crit}(\varpi)$ 的切空间由使它为零的 $X$ 组成。局部坐标中的简单计算表明，$\bar D^2\varpi_{x,I_{\mathcal E}X}=I_{\mathcal S}\bar D^2\varpi_{x,X}$。
\end{proof}

\SourcePage{227}
\StatementNumber{7}\begin{lemma}\label{st:15.7}
令 $x\in\operatorname{Crit}(\varpi)\setminus\operatorname{Crit}(\Psi)$，且 $\Psi(x)=w$。则 $x$ 是 $\varpi_w$ 的临界点；此外，交 $\operatorname{Crit}(\varpi)\cap\mathcal E_w$ 在 $x$ 处横截，当且仅当 $x$ 作为 $\varpi_w$ 的临界点非退化。
\end{lemma}
\begin{proof}
由于 $x$ 不是 $\Psi$ 的临界点，$D\varpi_x$ 的像与 $\mathcal S_w$ 的切空间横截。这蕴含 $\ker(D\varpi_x)=T(\mathcal E_w)_x$，且包含映射诱导的 $T(\mathcal S_w)_{\varpi(x)}\to\operatorname{coker}(D\varpi_x)$ 为同构。从而有 $D(\varpi_w)_x=0$，并且对 $X\in T(\mathcal E_w)_x$，可将\eref{15.13}与通常的二阶导数 $D^2(\varpi_w)_{x,X}$ 认同。若交 $\operatorname{Crit}(\varpi)\cap\mathcal E_w$ 横截，则对任何非零 $X\in T(\mathcal E_w)_x$，\eref{15.13}必非零，因而意味着该临界点非退化。同一论证也给出逆命题。
\end{proof}

\StatementNumber{8}\begin{lemma}\label{st:15.8}
在 $\Psi$ 的每个临界点 $x$ 处，映射 $\varpi|\operatorname{Crit}(\varpi):\operatorname{Crit}(\varpi)\to\mathcal S$ 为局部嵌入，而 $\Psi|\operatorname{Crit}(\varpi):\operatorname{Crit}(\varpi)\to\mathcal W$ 具有一个二阶分歧点。
\end{lemma}
\begin{proof}
由于 $D\varpi_x$ 秩为一，而 $D\Psi_x$ 为零，$D\psi_{\varpi(x)}$ 将 $\operatorname{coker}(D\varpi_x)$ 同构映到 $T\mathcal W_w$，其中 $w=\Psi(x)$。在此认同下，\eref{15.13}成为 $\Psi$ 在 $x$ 处的通常二阶导数；这意味着 $\operatorname{Crit}(\varpi)$ 的切空间在 $\CC$ 上由一个非零 $X$ 生成，满足对所有 $Y\in\ker(D\varpi_x)$ 有 $D^2\Psi_x(X,Y)=0$。已假设 $D^2\Psi_x$ 在 $\ker(D\varpi_x)$ 上非退化，故 $D\varpi_x(X)\ne0$，证明了第一项断言。非退化性还蕴含 $D^2\Psi_x(X,X)\ne0$，说明 $\Psi|\operatorname{Crit}(\varpi)$ 在 $x$ 处的二阶导数非零。
\end{proof}

\StatementNumber{9}\begin{lemma}\label{st:15.9}
在 $\Psi$ 的每个临界点及其在 $\mathcal S,\mathcal W$ 中的值附近，存在局部全纯坐标，使
\[
\varpi(x_1,\ldots,x_{n+2})=(x_1^2+x_2^2+\cdots+x_{n+2}^2,x_1),\qquad
\psi(y_1,y_2)=y_1.
\]
\end{lemma}
\begin{proof}
这是带参数全纯 Morse 引理的推论，该引理的证明与不带参数的情形相同。先在 $\mathcal S$ 与 $\mathcal W$ 上，以临界点的 $\varpi$ 值和 $\Psi$ 值为中心引入局部全纯坐标，使 $\psi$ 具有所需形式 $\psi(y_1,y_2)=y_1$；按假设，$I_{\mathcal S}$ 局部可积，故可做到。在这些坐标中局部写成 $\varpi=(\varpi_1,\varpi_2):\mathcal E\to\CC^2$，会发现 $D\varpi_2$ 在临界点非零，故水平集 $\{\varpi_2=0\}$ 光滑，且 $\varpi_1=\Psi$ 在此水平集上的限制在 $x$ 处具有非退化临界点。因此，附近的水平集 $\{\varpi_2=y_2\}$ 也光滑，$\varpi_1$ 限制到每一个水平集均有唯一非退化临界点，并全纯依赖于 $y_2$。于是可在每个水平集上找到全纯依赖于 $y_2$ 的局部坐标 $x_2,\ldots,x_{n+2}$，使 $\varpi_1$ 的限制具有通常的二次标准形。加上 $x_1=\varpi_2$，就得到 $\mathcal E$ 上一组坐标，使 $\varpi(x_1,\ldots,x_{n+2})=(f(x_1)+x_2^2+\cdots+x_{n+2}^2,x_1)$，
\SourcePage{228}
其中 $f$ 为某个全纯函数。由假设，$x=0$ 是 $\varpi_1=\Psi$ 的非退化临界点，故对某个 $c\ne0$，有 $f(t)=ct^2+O(t^3)$。最后，将 $x_1$ 换成 $f^{1/2}(x_1)$，并将 $y_2$ 换成 $f^{1/2}(y_2)$ 即可。
\end{proof}

现在讨论这些技术结果对于临界曲线及其像的几何意味着什么：
\SourceDiagram{diagram-15-critical}{0.43615338}
由于假设 $\operatorname{Critv}(\varpi)\cap\partial^h\mathcal S=\varnothing$，$\operatorname{Crit}(\varpi)$ 的边界等于它与 $\Psi^{-1}(\partial\mathcal W)$ 的交；又因 $\omega_{\mathcal E}$ 正合，$\operatorname{Crit}(\varpi)$ 的每个连通分支都必须有边界。这意味着没有分支能包含在 $\Psi$ 的某个纤维中，因此 $\Psi|\operatorname{Crit}(\varpi)$ 是有限分歧覆盖。由\stref{引理}{15.7}，对 $w\in\partial\mathcal W$ 要求 $\varpi_w$ 是 Lefschetz 纤维化，意味着在 $\partial\operatorname{Crit}(\varpi)$ 的某个邻域中，映射 $\varpi|\operatorname{Crit}(\varpi)$ 为嵌入，$\Psi|\operatorname{Crit}(\varpi)$ 为局部同构。应用一些基本伪全纯曲线理论，可见 $\varpi|\operatorname{Crit}(\varpi)$ 几乎处处单射，故其像 $\operatorname{Critv}(\varpi)$ 是没有多重覆盖分支的 $I_{\mathcal S}$ 全纯曲线。此外，曲线在边界 $\partial\operatorname{Critv}(\varpi)=\operatorname{Critv}(\varpi)\cap\partial^v\mathcal S$ 附近光滑，且 $\psi|\operatorname{Critv}(\varpi)$ 为局部同构。结合\stref{引理}{15.7}与\stref{引理}{15.8}可知，对 $w\in\operatorname{Critv}(\Psi)$，$\operatorname{Critv}(\varpi)$ 在 $\mathcal S_w$ 的某个邻域中也光滑，且每个这样的 $\mathcal S_w$ 包含 $\psi|\operatorname{Critv}(\varpi)$ 的唯一分歧点，该点还是二阶的。\sr{fig:15.1}{图~15.1}概括了这一情形。
\SourceFigure{figure-15-1}{0.53376866}{15.1}

定义\emph{伪临界值集} $\operatorname{Fakev}(\Psi,\varpi)$：它由那些 $w\in\mathcal W\setminus\operatorname{Critv}(\Psi)$ 组成，使 $\operatorname{Critv}(\varpi)$ 在 $\mathcal S_w$ 上具有奇点，或在 $\mathcal S_w$ 附近光滑但与其不横截相交。等价地，这些正是 $\mathcal E_w$ 光滑而\eref{15.12}不是 Lefschetz 纤维化的 $w$。由前面的讨论，$\operatorname{Fakev}(\Psi,\varpi)$ 是 $\operatorname{int}(\mathcal W)$ 的有限子集；特别地，$\mathcal W\setminus(\operatorname{Critv}(\Psi)\cup\operatorname{Fakev}(\Psi,\varpi))$ 仍连通。\znote{220}{伪临界值集是这一理论中的“例外集”：在 $w\notin\operatorname{Critv}(\Psi)$ 处 $\mathcal E_w$ 光滑，但 $\mathcal E_w\to\mathcal S_w$ 可能并非 Lefschetz 纤维化，因为 $\varpi_w$ 的临界点可能退化，或者 $\operatorname{Critv}(\varpi)$ 与 $\mathcal S_w$ 相切。前面关于临界曲线的讨论说明这样的 $w$ 只有有限多个：$\operatorname{Critv}(\varpi)$ 是无多重覆盖分支的全纯曲线，只在 $\Psi$ 的临界值处出现二阶分歧点。故从 $\mathcal W$ 中挖去有限多个点后，消失圈、匹配路径与单值变换的讨论全部可用；后文对一般 $w$ 陈述结论正是依据这一点。[来源：原书 §(15e)，印刷页 220。]}
